\section{Introduction}

The area-preserving H\'enon recurrence
\begin{equation}\label{eq:henon}
 x_{i+1}=a-x_i^2-x_{i-1}
\end{equation}
has exact periodic-orbit algebra, time reversal, and a natural symplectic
origin.  These features motivate its use as a structural test bed in the
foundational model of Wang~\cite{wang2026henon}.  The algebraic organization
of its chiral periodic orbits was developed by Endler and
Gallas~\cite{endler2006chiral}.

A preceding source audit and exact computation~\cite{hcs19} isolated an
apparent constant-placement error in a printed period-seven formula, adopted
the placement selected by an exact orbit fibre, and then proved the adopted
septic generically rather than trusting that fibre.  If
$a=\sigma^2-2\sigma$, the seven roots of the resulting
$P(\sigma,x)$ form one simple H\'enon cycle.  The scalar normalization $C$ has
genus three, but it forgets which neighbor is previous.  Its ordered edges
form a degree-14 cover carrying
\[
 \tau(x,y)=(a-x^2-y,x),\qquad
 J(x,y)=(x,a-x^2-y),
\]
with $\tau^7=J^2=1$ and $J\tau J=\tau^{-1}$.

The natural next question is not another scalar point count.  It is whether
the ordered-edge object is connected, what its compactified geometry is, and
whether H\'enon time produces arithmetic character sectors.  We answer all
three questions exactly.  The answer connects the specific H\'enon curve to
the classical theory of dihedral covers, double-coset endomorphisms, and real
multiplication developed by Ellenberg~\cite{ellenberg2001endomorphism} and
made explicit for genus-three curves by Hoffman, Liang, Sakai, and
Wang~\cite{hoffman2016genus3}.  Prym varieties of the relevant unramified
degree-seven covers have also been studied extensively~\cite{lange2016prym}.
We use these mechanisms; we do not claim them as new general theory.

The contributions for this H\'enon specialization are:
\begin{enumerate}
 \item the ordered-edge curve $E$ is the connected $D_7$ splitting curve of
 $P$, with $g(E)=8$;
 \item the complete quotient diagram is explicit: the rotation quotient is
 $B:w^2=Q_6(\sigma)$ of genus two, the reflection quotient is $C$ of genus
 three, and $E\to B$ is unramified cyclic of degree seven;
 \item the quadratic extension $E\to C$ is
 $\Q(C)(\sqrt{Q_6})/\Q(C)$ and its six branch points have an exact coordinate
 formula;
 \item $\Jac(E)\sim_\Q\Jac(B)\times\Jac(C)^2$;
 \item one H\'enon step induces a Rosati-self-adjoint correspondence on
 $\Jac(C)$ generating $\Q(\zeta_7+\zeta_7^{-1})$; and
 \item at $p=5,11,13$, tame vertical-inertia exclusion and a two-chart
 normalization comparison certify the scalar and ordered-edge local factors.
\end{enumerate}

The last item is the relevant spectral structure.  It explains why a
degree-six scalar Frobenius polynomial should be the norm of one quadratic
over a real cubic field, and the three certified rows satisfy that law
exactly.  It also imposes a limitation: forgetting the $\tau$ labels, the
ordinary cohomology of $E$ contains only $H^1(B)$ and two copies of
$H^1(C)$.  Orientation refines the spectrum equivariantly but does not create
new ordinary eigenvalues.
